\section{Discussion}
\label{sec:discussion}

\paragraph{What the protocol caught.}
The apparent increment against a single recalibrated HAR---38 of 69 cells on the seed-ensemble primary---decomposes into exactly the two artefacts the protocol targets. Part is baseline weakness: against the maximal price pool, 9 of 69 cells survive Holm. Part is firm identity: against the identity-augmented reference, 8 of 69 survive; a \emph{zero-text} firm-mean forecast reproduces the apparent gain, beating the plain reference in 53 of 69 cells, and 40 of 45 long-form cells flip negative---text \emph{hurts} once the reference knows the firm. The two survivor sets are disjoint---their intersection is empty (0 of 69 under Holm)---and the strictest gate, placebo-confirmed genuine cells clearing both controls, leaves nothing. This is the quantitative anatomy of an inflated ``representation helps'' claim: each control alone leaves a remnant, and the other control explains it.

\paragraph{Detectable, attributable, bankable.}
Encompassing regressions \citep{harvey1998encompassing} return $g>0$ in 48 of 48 tests (HAC $t$ from $+2.3$ to $+21.1$): conditional on HAR, text forecasts are \emph{detectable}, so the near-null is not ``text is noise''. But detectability is not attribution: the identity control shows that much of what $g$ detects is \emph{who filed}, content a stronger reference already carries. And attribution is not realisability: under the deployable expanding combiner most cells collapse (event-driven TF--IDF at $h{=}5$ falls from 8 of 16 significant quarters to 1 of 16), and the economics are nil (median $\Delta$Sharpe $+0.003$; significant in 1 of 18 cells, about the multiplicity floor). Conflating the three claims is how weak-baseline studies overclaim.

\paragraph{Why periodic filings approximate a firm dummy.}
Our models read filings as \emph{levels}, not changes. Periodic 10-K/10-Q language is firm-stable---the risk-factor and MD\&A sections that carry the long-form signal (Section~\ref{sec:ablations}) recur from one filing to the next---so periodic long-form filings approximate a firm dummy, and the identity control makes that explicit. Event-driven 8-Ks exist because something happened; filing-specific news lives there, and that is exactly where the residual sits: $+$0.2--0.5\% QLIKE in a handful of 8-K cells, cross-sectional (which-firm) in 33 of the 38 genuine cells rather than regime-timing.

\paragraph{Prompting versus pooling.}
On byte-identical curated excerpts, the prompted 32B LLM adds a genuine long-form increment ($+1.79/+2.25/+0.27\%$ at $h{=}5/10/20$) while pooled same-lineage embeddings add none (0 of 3): the dissociation is mechanism, not curation. Contamination controls bound the reading: a date-plus-ticker prompt with \emph{zero} text reproduces 31--77\% of the full-text increment in the well-identified cells, though full text still adds beyond that identity prompt (6 of 6 Holm) and survives an era split at the model's training cutoff (5 of 6 under Holm). The increment also does not replicate across families---Yi-1.5-34B and Phi-4 yield no significant positive cell anywhere, though both are smaller, older, and mode-collapse---so we treat the dissociation as bounded evidence about elicitation, not a bankable capability.

\paragraph{Fit without generalisation.}
The neural text models are not undertrained. They fit the training period with $R^2$ of 0.30--0.65, well above HAR's 0.17--0.31, yet their test $R^2$ turns negative; training-free controls---ridge and dictionary models with no epochs to tune---fail identically. The fair-compute objection therefore has no purchase: the fixed recipe is deliberate, and the signature is overfitting to non-transferable features, not insufficient optimisation.

\paragraph{Guidance for benchmark builders.}
Five practices are cheap and portable: recalibrate the reference before crediting a challenger; add an entity-identity term to the reference; cluster inference on the unit that shares shocks (here, the trading day); gate significance with placebos; and report the \emph{reference interval}, not a single point. We scope the lesson honestly: we demonstrate the artefact for firm identity on one benchmark, not as a general theorem about learned representations.

\paragraph{What would change the verdict.}
Each condition of the design points to an extension that raises or relocates the bar. Intraday realised measures \citep{andersen2003modeling} or option-implied references \citep{blair2001forecasting} would strengthen the price side. Other markets and small caps---less analysed, more slowly impounded---relocate it, and are where disclosure text has its best remaining chance. Richer conditioning on event type could localise value: the one cell that survives all four identity specifications (event-driven, prompted LLM, $h{=}5$) is where we would look first. In every direction, the burden stays on text to clear an identity-aware, recalibrated reference.
